%! Parametric Functions and Rates of Change — Unit 8, Lesson 8.2 — Group Activity (Answer Key)
\documentclass[10pt]{article}
\usepackage{precalculus-article}
\usepackage{precalculus-key}   % pulls in -boxes; do NOT also load -boxes
\newcommand{\TallMath}[1]{$\displaystyle #1\rule[-1.4em]{0pt}{3.2em}$}

\begin{document}

\pageheader{Unit 8, Lesson 8.2}{Group Activity --- Answer Key}
\namepartnerperiod

\vspace{0.12in}

\begin{scenariobox}[Two Robots, One Path]{plumlight}
Two robots traverse the same parabolic path in a warehouse.
\[
  \text{Robot A:}\quad g_A(t) = (t,\; t^2), \quad 0 \le t \le 3
  \qquad\qquad
  \text{Robot B:}\quad g_B(t) = (3-t,\; (3-t)^2), \quad 0 \le t \le 3
\]
Robot A starts at $(0, 0)$ and Robot B starts at $(3, 9)$.
\end{scenariobox}

\vspace{0.10in}

\begin{notesbox}{Part 1 --- Tables and Shared Path}

\textbf{1.}\quad Complete the table for each robot.

\vspace{6pt}
\begin{center}
\renewcommand{\arraystretch}{1.7}
\begin{tabular}{|c|c|c|}
\hline
$t$ & Robot A position $(x, y)$ for $g_A(t) = (t,\; t^2)$ &
      Robot B position $(x, y)$ for $g_B(t) = (3-t,\; (3-t)^2)$ \\
\hline
$0$ & \ans{$(0,\,0)$} & \ans{$(3,\,9)$} \\
\hline
$1$ & \ans{$(1,\,1)$} & \ans{$(2,\,4)$} \\
\hline
$2$ & \ans{$(2,\,4)$} & \ans{$(1,\,1)$} \\
\hline
$3$ & \ans{$(3,\,9)$} & \ans{$(0,\,0)$} \\
\hline
\end{tabular}
\end{center}

\vspace{8pt}
\textbf{2.}\quad Do both robots visit the same set of ordered pairs?

\vspace{4pt}
\ans{Yes. Both robots visit $\{(0,0),\,(1,1),\,(2,4),\,(3,9)\}$ --- the same four points, just in opposite order. Robot A goes from $(0,0)$ to $(3,9)$; Robot B goes from $(3,9)$ to $(0,0)$.}

\end{notesbox}

\vspace{0.08in}

\begin{notesbox}{Part 2 --- Direction of Motion}

\textbf{3.}\quad At $t = 1$, is Robot A moving right or left? Up or down?

\vspace{4pt}
Compare $x_A(0) = 0$ and $x_A(2) = 2$: $x$ is \ans{increasing} $\Rightarrow$ Robot A moves \ans{right}.

Compare $y_A(0) = 0$ and $y_A(2) = 4$: $y$ is \ans{increasing} $\Rightarrow$ Robot A moves \ans{up}.

\vspace{0.10in}
\textbf{4.}\quad At $t = 1$, is Robot B moving right or left? Up or down?

\vspace{4pt}
Compare $x_B(0) = 3$ and $x_B(2) = 1$: $x$ is \ans{decreasing} $\Rightarrow$ Robot B moves \ans{left}.

Compare $y_B(0) = 9$ and $y_B(2) = 1$: $y$ is \ans{decreasing} $\Rightarrow$ Robot B moves \ans{down}.

\vspace{0.10in}
\textbf{5.}\quad Do the two robots move in the same direction at $t = 1$?

\vspace{4pt}
\ans{No. At $t = 1$, Robot A moves right and up, while Robot B moves left and down. This illustrates EK 4.3.A.3: the same curve can be traversed in opposite directions by different parametrizations.}

\begin{teachernote}
  Students sometimes assume both robots must move the same direction since they're on the
  same path. Q5 is the key conceptual moment --- help groups articulate that direction is
  a property of the parametrization, not the curve itself.
\end{teachernote}

\end{notesbox}

\vspace{0.08in}

\begin{notesbox}{Part 3 --- Slope of the Curve}

\textbf{6.}\quad Using Robot A's data, slope between $t = 0$ and $t = 2$.

\vspace{4pt}
$x_A(0) = $ \ans{$0$} \qquad $x_A(2) = $ \ans{$2$} \qquad $\Delta x = $ \ans{$2$}

$y_A(0) = $ \ans{$0$} \qquad $y_A(2) = $ \ans{$4$} \qquad $\Delta y = $ \ans{$4$}

\vspace{4pt}
Slope $= \dfrac{\Delta y}{\Delta x} = \dfrac{4}{2} = $ \ans{$2$}

\vspace{10pt}
\textbf{7.}\quad Using Robot B's data for the same two $(x,y)$ locations.

\vspace{4pt}
Robot B passes through $(0, 0)$ at $t = $ \ans{$3$}
and through $(2, 4)$ at $t = $ \ans{$1$}.

$\Delta x = x_B(1) - x_B(3) = 2 - 0 = $ \ans{$2$} \qquad
$\Delta y = y_B(1) - y_B(3) = 4 - 0 = $ \ans{$4$}

\vspace{4pt}
Slope $= \dfrac{\Delta y}{\Delta x} = $ \ans{$2$}

\vspace{4pt}
\ans{Both give slope $= 2$. This makes sense because slope is a property of the \emph{curve} between two geometric points, not of the parametrization. Different parametrizations of the same path must give the same secant slope.}

\end{notesbox}

\vspace{0.08in}

\begin{notesbox}{Part 4 --- When Does Each Robot Arrive?}

\textbf{8.}\quad When does each robot pass through $(1, 1)$?

\vspace{4pt}
Robot A: Set $x_A(t) = t = 1 \Rightarrow$ \ans{$t = 1$}. Check: $y_A(1) = 1$ \checkmark

Robot B: Set $x_B(t) = 3 - t = 1 \Rightarrow t = 2$. \ans{$t = 2$}. Check: $y_B(2) = 1$ \checkmark

\vspace{0.10in}
\textbf{9.}\quad When does each robot pass through $(0, 0)$? Through $(3, 9)$?

\vspace{4pt}
$(0, 0)$: Robot A at \ans{$t = 0$}; Robot B at \ans{$t = 3$}.

$(3, 9)$: Robot A at \ans{$t = 3$}; Robot B at \ans{$t = 0$}.

\begin{teachernote}
  Part 4 reinforces EK 4.3.A.2: at the same geometric point, different parametrizations
  arrive at different values of $t$. Ask groups: ``Does the slope of the secant change
  depending on which robot you use to compute it?'' (No.)
\end{teachernote}

\end{notesbox}

\end{document}
